\subsection{向量空间的标量扩张}

回顾（I，§ 9，第 1 号，定理 2），域 K 到非零环 A 的同态必然是单射。

命题 18. 设 \(\rho\) 为域 \(K\) 到环 \(A\) 中的一个同态。对于每一个由向量 \(K\) -空间和 \(K\) -线性映射组成的正合序列

\[
\mathrm{E} ^ {\prime} \xrightarrow {u} \mathrm{E} \xrightarrow {v} \mathrm{E} ^ {\prime \prime}
\]

该序列

\[
\mathrm{E} _ {(\mathrm{A})} ^ {\prime} \xrightarrow {u _ {(\mathrm{A})}} \mathrm{E} _ {(\mathrm{A})} \xrightarrow {v _ {(\mathrm{A})}} \mathrm{E} _ {(\mathrm{A})} ^ {\prime \prime}
\]

是正合的。

这是第 7 号命题 14 的一个特例，需考虑到 § 1 第 4 号中的注记 4。

推论。对于每一个K-线性映射 \(f: \mathrm{E}' \to \mathrm{E}\) , \(\operatorname{Im}(f_{(\mathbf{A})}) = (\operatorname{Im}(f))_{(\mathbf{A})}\) , \(\operatorname{Ker}(f_{(\mathbf{A})}) = (\operatorname{Ker}(f))_{(\mathbf{A})}\) , \(\operatorname{Coker}(f_{(\mathbf{A})}) = (\operatorname{Coker}(f))_{(\mathbf{A})}\) ，在典范同构的意义下。

命题 19. 设 \(\rho\) 为域 K 到环 A 中的一个单同态。对于 K 上的每个左向量空间 E，典范映射 \(\phi: \mathbf{E} \to \rho^{*}(\mathbf{E}) = \mathbf{A} \otimes_{\mathbb{K}} \mathbf{E}\) 是单射。此外，对 E 的每个向量子空间 \(\mathbf{E}'\) ， \(\rho^{*}(\mathbf{E}') = \mathbf{A} \otimes_{\mathbb{K}} \mathbf{E}'\) 典范地等同于 \(\mathbf{A} \otimes_{\mathbb{K}} \mathbf{E}\) 的一个直因子子 A-模，并且在这种等同下，

\[
(\mathbf {A} \otimes_ {\mathbf {K}} \mathbf {E} ^ {\prime}) \cap \phi (\mathbf {E}) = \phi (\mathbf {E} ^ {\prime}).\tag{35}
\]

第一个断言是 § 5 第 1 号命题 4 的特例；第二个是第 7 号命题 14 的特例；最后，为证明 (35)，只需在 A（视为右 K-模）中取一组基 \((a_{\lambda})_{\lambda \in L}\) ，使得 \(a_{\lambda_0} = 1\) 对某个指标 \(\lambda_0\) 成立（第 1 号，定理 2）；A \(\otimes_{\mathbb{K}}\) E 的元素可以唯一地写成 \(\sum_{\lambda} a_{\lambda} \otimes x_{\lambda}\) ，其中 \(x_{\lambda} \in \mathbf{E}\) ，而这样的元素属于 A \(\otimes_{\mathbb{K}}\) E' 的必要且充分条件是： \(x_{\lambda} \in \mathbf{E}'\) 对所有 \(\lambda\) 成立。另一方面， \(\phi(\mathbf{E})\) 的元素是那些使得 \(x_{\lambda} = 0\) 对 \(\lambda \neq \lambda_0\) 成立的元素；对元素 \(\sum_{\lambda} a_{\lambda} \otimes x_{\lambda}\) 属于 (A \(\otimes_{\mathbb{K}}\) E') \(\cap\) \(\phi(\mathbf{E})\) ，必要且充分的条件是 \(x_{\lambda} = 0\) 对 \(\lambda \neq \lambda_0\) 成立，且 \(x_{\lambda_0} \in \mathbf{E}'\) ，由此得出结论。

推论。设 \(\rho\) 为域 K 到环 A 中的一个单同态。K-线性映射 \(f: \mathbf{E} \to \mathbf{F}\) （其中 E 和 F 是 K 上的两个向量空间）为单射（相应地为满射、零映射），当且仅当 \(f_{(\mathbf{A})}: \mathbf{E}_{(\mathbf{A})} \to \mathbf{F}_{(\mathbf{A})}\) 为单射（相应地为满射、零映射）。

这直接由命题 19 及命题 18 的推论得出。

命题 20. 设 \(\rho\) 为域 \(\mathbf{K}\) 到环 \(\mathbf{A}\) 中的一个同态。对于每一个左向量空间 \(\mathbf{E}\) （ \(\mathbf{K}\) 上的），典范右 A-模同态

\[
\upsilon : \left(\mathrm{E} ^ {*}\right) _ {(\mathrm{A})} \rightarrow \left(\mathrm{E} _ {(\mathrm{A})}\right) ^ {*}
\]

（§ 5，第 4 号）是单射；当 E 为有限维时，它是双射。

第二个断言由 § 5 第 4 号命题 8 得出。为证明第一个断言，我们注意到 \((\mathbf{E}^{*})_{(\mathbf{A})}\) 的每个元素都可以写成 \(\sum_{i} x_{i}^{*} \otimes \alpha_{i}\) ，其中 \(\alpha_{i} \in \mathbf{A}\) ，而 \((x_{i}^{*})_{1 \leqslant i \leqslant n}\) 是 \(\mathbf{E}^{*}\) 中的一个自由族；在 \((\mathbf{E}_{(\mathbf{A})})^{*}\) 中对应于它的是线性形式 \(y^{*}\) ，使得 \(y^{*}(1 \otimes x) = \sum_{i} \rho(\langle x, x_{i}^{*} \rangle)\alpha_{i}\) 对所有 \(x \in \mathbf{E}\) 成立。但是，E 中存在一个族 \((x_{i})_{1 \leqslant i \leqslant n}\) ，使得 \(\langle x_{i}, x_{j}^{*} \rangle = \delta_{ij}\) （第 5 号，定理 7 的推论 2），由此 \(y^{*}(1 \otimes x_{i}) = \alpha_{i}\) ；因此关系 \(y^{*} = 0\) 蕴含 \(\alpha_{i} = 0\) 对所有 \(i\) 成立，这证明了我们的断言。

命题 21. 设 K 是一个域，L 是 K 的一个扩域。

\begin{enumerate}
\def\labelenumi{(\roman{enumi})}
\item
  对于每个向量空间 \(\mathbf{E}\) （ \(\mathbf{K}\) 上的）， \(\dim_{\mathbf{L}}(\mathrm{E}_{(\mathbf{L})}) = \dim_{\mathbf{K}}\mathbf{E}\) .
\item
  对于每个 K-线性映射 \(u: \mathbf{E} \to \mathbf{F}\) （其中 \(\mathbf{E}\) 和 \(\mathbf{F}\) 是 \(\mathbf{K}\) 上的向量空间）， \(\operatorname{rg}(u_{(L)}) = \operatorname{rg}(u)\) .
\end{enumerate}

若 \((e_{\iota})_{\iota\in I}\) 是 \(E\) 在 \(K\) 上的一组基，则 \((1\otimes e_{\iota})_{\iota\in I}\) 是 \(\mathbf{E}_{(L)}\) 在 \(L\) 上的一组基（§ 5，第 1 号，命题 4），由此得第一个断言；第二个断言由第一个断言及以下事实推出： \(u_{(L)}(\mathbf{E}_{(L)})\) 依命题 18 的推论典范地等同于 \((u(\mathbf{E}))_{(L)}\) 。

命题 22. 设 K 是一个交换域， \(\rho: \mathrm{K} \to \mathrm{A}\) 为一个单射的中心同态，且 E、F 是 K 上的两个向量空间。则典范同态

\[
\omega : \mathrm{A} \otimes_ {\mathrm{K}} \operatorname{Hom} (\mathrm{E}, \mathrm{F}) \rightarrow \operatorname{Hom} _ {\mathrm{A}} \left(\mathrm{E} _ {(\mathrm{A})}, \mathrm{F} _ {(\mathrm{A})}\right)\tag{36}
\]

（§ 5，第 3 号，公式 (17)）是单射的；如果 A 或 E 是 K 上的有限维向量空间，则它是双射的。

这是 § 5 第 3 号命题 7 的一个特例。

